\begin{resumo}[Abstract]

The search for predicting asset market behavior has motivated the development of economic science and different analytical approaches, such as fundamental and technical analysis. However, market movements are a direct reflection of human emotions and behavior. Previously, accessing this information was restricted to well-positioned professionals with a large network of contacts. However, the use of Artificial Intelligence has been dominating various areas that benefit from large amounts of available data. In order to benefit from this technological tool in the financial market, this work proposes the use of advanced text data processing techniques, such as Natural Language Processing (NLP), to extract information from various internet sources. The extracted information will be used as input data for a machine learning model based on Convolutional Neural Networks (CNN) to classify them on a scale of 0 to 100, quantifying market sentiment towards a specific asset. This sentiment score will be subsequently used as a weighting parameter in a final asset allocation model for an investment portfolio. The implementation of this tool aims to contribute with reliable information for decision-making, considering the defined risk-return profile.

\textbf{Keywords}: Convolutional Neural Networks. Natural Language Processing. Web Scraping. Market sentiment. Investment portfolio.

{\color{red}
Resumo traduzido para outros idiomas, neste caso, inglês. Segue o formato do resumo feito na língua vernácula. As palavras-chave traduzidas, versão em língua estrangeira, são colocadas abaixo do texto precedidas pela expressão Keywords, separadas por ponto.

\vspace{\onelineskip}

\textbf{Keywords}: Word. Word. Word. 
}
	
\end{resumo}
